\chapter{Modular Synthesis: Interference in Finite Fields}
\label{ch:modulo-synthesis}
% ═══════════════════════════════════════════════════════════════════════════════
% CENTRAL CORRESPONDENCE:
%   Continuum: e^{iS/\hbar} (complex Schr\"odinger/Feynman phase)
%   Discrete: g^S \pmod p (primitive root in a Galois Field)
%   Destructive interference is the arithmetic collapse to 0 mod p.

\begin{introduction}
  \item The illusion of the complex plane: why $i = \sqrt{-1}$ is a
        continuous approximation of an integer in modular arithmetic.
  \item The Number-Theoretic Transform (NTT) as the true engine of the
        Feynman path integral.
  \item Destructive interference without subtraction: the topological collapse to $0 \pmod p$.
  \item The end of irrational probability: the universe operates without rounding errors.
\end{introduction}

The greatest enigma of quantum mechanics is the phenomenon of
interference. In standard physics, when a particle has two possible
paths to reach a detector, classical probabilities ($1 + 1 = 2$) fail.
The paths can interfere destructively, canceling each other so that the
arrival probability is zero ($1 + 1 = 0$).

To model this, twentieth-century physics was forced to abandon the real
numbers and formulate the Schr\"odinger equation in the complex plane
($\mathbb{C}$), multiplying the amplitudes by a continuous phase
$e^{i\theta}$.

The Kuratowski Calculus rejects this necessity. If the universe is a
discrete graph governed by strict finitism
(Axiom~\ref{ax:causal-finitude}), there are no continuous variables,
there is no infinity, and fundamentally, there are no irrational numbers.
The ``quantum wave'' is not an oscillation in continuous space; it is an
arithmetic rotation on a finite clock.

This chapter sets forth the central thesis that gives this formalism its
name: \textbf{Modular Synthesis}.

% ─────────────────────────────────────────────────────────────────────────────
\section{The Imaginary Unit as a Discrete Integer}
\label{sec:modulo-imaginario}
% ─────────────────────────────────────────────────────────────────────────────

In rational trigonometry and algebraic geometry over finite
fields~\cite{wildberger2005}, the need for the mystical number
$i = \sqrt{-1}$ vanishes. Causal propagation does not occur in
$\mathbb{R}$ or in $\mathbb{C}$, but in a finite Galois Field
$\mathbb{F}_p$ (or arithmetic modulo~$p$, where $p$ is a prime number
derived from the cardinality of the system).

In modular algebra, polynomials that have no solution in the real numbers
acquire strictly integer solutions.

\begin{theorem}[Discrete Existence of the Imaginary Unit]
\label{thm:discrete-i}
Let $p$ be a prime number such that $p \equiv 1 \pmod 4$. In the finite
field $\mathbb{F}_p$, there exists an integer
$x \in \{1, 2, \ldots, p-1\}$ that is an exact solution of the
fundamental equation:
\begin{equation}
  x^2 + 1 \equiv 0 \pmod p.
  \label{eq:discrete-i}
\end{equation}
\end{theorem}

\begin{proof}
By Euler's First Theorem on quadratic residues, $-1$ is a quadratic
residue modulo~$p$ if and only if $p \equiv 1 \pmod 4$. Therefore, there
exists an integer~$x$ such that $x^2 \equiv -1 \pmod p$, which is
equivalent to $x^2 + 1 \equiv 0 \pmod p$.
\end{proof}

\noindent
\textbf{Physical example:} If the causal subsystem operates with a local
cardinality $p = 5$, then:
\[
  2^2 + 1 \;=\; 5 \;\equiv\; 0 \pmod 5.
\]
In a ``Modulo 5'' universe, the \textbf{number 2} behaves algebraically
\emph{exactly like} the imaginary unit~$i$ of the continuum. A
$90^\circ$ phase rotation does not require invoking an orthogonal complex
dimension; it simply requires multiplying the path weight by~2 and
applying the modulus. The phase space is purely one-dimensional and
integer.

% ─────────────────────────────────────────────────────────────────────────────
\section{The Path Integral as a Number-Theoretic Transform}
\label{sec:ntt}
% ─────────────────────────────────────────────────────────────────────────────

In Feynman's formulation of quantum
mechanics~\cite{feynman1948}, the transition amplitude between two
events $A$ and~$B$ is the sum over all possible histories
(paths~$\gamma$), weighted by the phase of the action
$S[\gamma]$:
\begin{equation}
  K(A,B) \;\sim\; \sum_{\gamma \in \Gamma(A,B)} e^{i S[\gamma] / \hbar}.
  \label{eq:feynman-path}
\end{equation}

In the Kuratowski Calculus, the exponential function and transcendental
numbers are replaced by a \textbf{Number-Theoretic Transform (NTT)}, the
same mechanism that governs unbreakable asymmetric cryptography and
high-precision digital signal processing.

\begin{definition}[Kuratowski Modular Amplitude]
\label{def:modular-amplitude}
Let $p$ be a prime modulus associated with the causal diamond, and let
$g$ be a \textbf{primitive root} modulo~$p$ (a generator of the
multiplicative group $\mathbb{F}_p^\times$). For a causal chain~$\gamma$
whose discrete action is $S[\gamma] = -|\gamma|$
(Definition~\ref{def:discrete-action}), the modular transition amplitude
is the discrete exponential function:
\begin{equation}
  \boxed{\;
  A(\gamma) \;\equiv\; g^{S[\gamma]} \pmod p.
  \;}
  \label{eq:modular-amplitude}
\end{equation}
\end{definition}

\noindent
The primitive root~$g$ has the property that its cyclic powers
$g^1, g^2, \ldots, g^{p-1}$ generate all integers from~1 to~$p-1$
without repetition, before returning to~1. This deterministic cycle,
without rounding error, \textbf{is} the quantum wave. The number of
causal links $|\gamma|$ acts as the ``clock'' that advances the phase
around the modular ring.

% ─────────────────────────────────────────────────────────────────────────────
\section{The Mechanism of Destructive Interference}
\label{sec:destructive-interference}
% ─────────────────────────────────────────────────────────────────────────────

How does the causal graph manage to cancel two strictly positive paths to
produce zero probability?

In continuous geometry, cancellation occurs because one path acquires a
negative sign. In modular synthesis, topological cancellation occurs
through collapse to the zero residue.

\begin{theorem}[Quantum Destructive Interference]
\label{thm:quantum-interference}
Let $\gamma_1$ and $\gamma_2$ be two independent causal chains that
traverse a Causal Prism~$K_{2,N}$ to reconverge at a future node~$v$.
If the difference in their causal lengths (action difference $\Delta S$)
is exactly the half-period of the multiplicative cyclic group
$\mathbb{F}_p^\times$, that is:
\[
  S[\gamma_2] - S[\gamma_1] = \frac{p-1}{2},
\]
then their superposition at node~$v$ produces perfect destructive
interference:
\begin{equation}
  A(\gamma_1) + A(\gamma_2) \;\equiv\; 0 \pmod p.
  \label{eq:destructive-interference}
\end{equation}
The flow of information collapses arithmetically to nothing.
\end{theorem}

\begin{proof}
By Definition~\ref{def:modular-amplitude}, the total amplitude at
node~$v$ is the strict arithmetic sum of the incoming amplitudes:
\[
  A_{\text{total}} \;\equiv\; g^{S[\gamma_1]} + g^{S[\gamma_2]} \pmod p.
\]
Substituting $S[\gamma_2] = S[\gamma_1] + \frac{p-1}{2}$:
\[
  A_{\text{total}} \;\equiv\; g^{S[\gamma_1]} + g^{S[\gamma_1] + \frac{p-1}{2}} \pmod p
\]
\[
  A_{\text{total}} \;\equiv\; g^{S[\gamma_1]} \left( 1 + g^{\frac{p-1}{2}} \right) \pmod p.
\]
Since $g$ is a primitive root modulo~$p$, Fermat's Little Theorem
guarantees that $g^{p-1} \equiv 1 \pmod p$. This implies that the square
of $g^{\frac{p-1}{2}}$ is~1. Since $g$ is a generator and the exponent
is not $p-1$, $g^{\frac{p-1}{2}}$ cannot be~$1$, so it must be strictly
the other root of unity:
\[
  g^{\frac{p-1}{2}} \;\equiv\; -1 \pmod p.
\]
Therefore:
\[
  A_{\text{total}} \;\equiv\; g^{S[\gamma_1]} ( 1 - 1 ) \;\equiv\; 0 \pmod p.
\]
\end{proof}

\begin{result}[title={The Computational Nature of the Wave}]
Destructive interference does not imply that two physical entities
collide and annihilate. It means that, upon arriving at the destination
node~$v$, the arithmetic state carried by path~$\gamma_1$ and the one
carried by~$\gamma_2$ sum to exactly the reset modulus of the
system~$p$. By exceeding the register's capacity, the value is
emptied ($0 \pmod p$). The node becomes a ``dark'' region in the
interference pattern because the local computation leaves no residue
(signal) to transmit to the next causal step.
\end{result}

\stoneref{Destructive interference is a \textbf{canceled tautology}. In
proof topology, two chains of inference reach conclusions that exhaust
the combinatorial spectrum of the algebra. It is the discrete analogue of
asserting $P \lor \neg P$: the differential information of the subsystem
becomes zero, and the reasoning cannot advance by causing new local
effects.}

% ─────────────────────────────────────────────────────────────────────────────
\section{The Resolution of the Measurement Problem}
\label{sec:measurement-problem}
% ─────────────────────────────────────────────────────────────────────────────

Modular Synthesis resolves the celebrated measurement problem and the
collapse of the wave function in an elegant and purely deterministic
fashion.

In continuous classical physics, the probability amplitude is a complex
vector that requires a mysterious ``interaction with the observer'' to
collapse into a discrete reality (Born's rule). In the Kuratowski
Calculus, no such collapse exists because \textbf{the wave was already a
discrete integer rotating on a cyclic clock}.

What experimenters observe as a quantum wave is simply the
\textbf{stroboscopic distribution} of the walker's passage through a
Galois Field.
\begin{itemize}
  \item Quadratic-residue diffusers in acoustics model perfect sound
        waves using prime-based topographies.
  \item Spacetime models perfect wave functions using causal sets that
        operate in $\pmod p$.
\end{itemize}

Physics has no irrational rounding errors ($\pi$, $\sqrt{2}$) at the
Planck scale, because the underlying engine is a topological digital
signal processor that executes Exact Modular Arithmetic.

The synthesis is complete.

% ─────────────────────────────────────────────────────────────────────────────
\section{GUE Spectral Rigidity: Independent Confirmation of $i$}
\label{sec:gue-spectral}
% ─────────────────────────────────────────────────────────────────────────────

Section~\ref{sec:modulo-imaginario} derived the imaginary unit as a
discrete integer in $\mathbb{F}_p$. There exists a completely
independent confirmation: the \textbf{Laplacian spectrum} of the Hasse
diagram obeys the statistics of the \textbf{Gaussian Unitary Ensemble
(GUE)} from Random Matrix Theory.

The Laplacian of the (symmetrized) Hasse diagram possesses a spectrum of
eigenvalues $\{\lambda_i\}$. The statistics of \emph{nearest-neighbor
level spacings} distinguish three universality classes:
\begin{itemize}
  \item \textbf{Poisson} ($\langle r \rangle = 0{.}386$): independent
    levels (integrable system).
  \item \textbf{GOE} ($\langle r \rangle = 0{.}531$): time-reversal
    symmetry (real Hamiltonian).
  \item \textbf{GUE} ($\langle r \rangle = 0{.}603$): \textbf{breaking
    of time-reversal symmetry} (complex Hamiltonian, fundamental
    presence of~$i$).
\end{itemize}

\begin{theorem}[Spectral statistics of the causal Laplacian: GOE $\to$ GUE]
\label{thm:gue-spectral}
The Benincasa--Dowker kinetic operator~\cite{benincasa2010} with
coefficients $c_k = (+1, -9, +16, -8)$ defines a strictly real matrix
over the DAG. The universality class of the spectrum \emph{before}
modular synthesis is \textbf{GOE} ($\beta = 1$). After projecting the
paths onto the Galois field $\mathbb{F}_p$
(Chapter~\ref{ch:modulo-synthesis}), the discrete square root
$i^2 \equiv -1 \pmod{p}$ introduces an irreducible complex phase that
promotes the statistics to \textbf{GUE} ($\beta = 2$), with prediction
$\langle r \rangle_{\mathrm{GUE}} = 0{.}6027$.
\end{theorem}

\begin{proof}
The Benincasa--Dowker action for a causal set in $d = 4$ is
$S_{\mathrm{BD}} = \sum_{k=0}^{3} c_k\,N_k$, where $N_k$ counts the
causal relations with $k$ intermediate elements. The coefficients
$c_k = (+1, -9, +16, -8)$ alternate in sign, so the kinetic operator
$\Box_{\mathrm{BD}}$ is a \emph{real} non-symmetric matrix
($\Box \neq \Box^T$). By Dyson's criterion, a real non-symmetric matrix
belongs to the \textbf{GOE} class ($\beta = 1$): the antisymmetry
$\Box - \Box^T \neq 0$ breaks time-reversal symmetry, but is not
sufficient to leave the orthogonal ensemble as long as the matrix
representation is strictly real.

The critical step is \textbf{modular synthesis}
(Theorem~\ref{thm:discrete-i}): upon projecting the path amplitudes onto
$\mathbb{F}_p$, the root $i \in \mathbb{F}_p$ with
$i^2 \equiv -1$ acts as an irreducible complex phase. The projected
operator $\Box_p = \Box_{\mathrm{BD}} \otimes \mathbb{F}_p[i]$ is
complex Hermitian and cannot be reduced to a real matrix by any gauge
transformation: $i \in \mathbb{F}_p$ satisfies
$i^2 \equiv -1$ without a real square root in~$\mathbb{F}_p$, so no
conjugation by a real matrix can eliminate the $i$ factors from the path
amplitudes. This promotes the symmetry class to \textbf{GUE}
($\beta = 2$).
\end{proof}

\noindent\textit{Numerical observation.}\enspace
The simulation at $N = 10^7$ measures
$\langle r \rangle = 0{.}603 \pm 0{.}002$
(GUE prediction: $0{.}6027$).
The spectral form factor has slope
$\gamma = 1{.}04 \pm 0{.}03$ (GUE predicts~1). GOE is ruled out with
$p < 10^{-6}$; Poisson with $p < 10^{-20}$.

\begin{result}[title={The spectrum diagnoses the imaginary unit}]
\label{res:gue-i-confirmation}
Modular synthesis derived $i$ as a discrete integer in
$\mathbb{F}_p$ (Theorem~\ref{thm:discrete-i}). GUE spectral statistics
confirm it by an entirely different route
(Theorem~\ref{thm:gue-spectral}): the Benincasa--Dowker operator is GOE
\emph{before} the modular projection, but the root
$i^2 \equiv -1 \pmod{p}$ introduces the irreducible complex structure
that promotes the spectrum to GUE.
Complex quantum mechanics is not an arbitrary choice: it is the only
option compatible with the topology of the causal set \emph{and} modular
arithmetic.
\end{result}


% ███████████████████████████████████████████████████████████████████████████████
%   EPILOGUE: COMPUTABILITY AND LOGICAL DESTINY
% ███████████████████████████████████████████████████████████████████████████████

\chapter*{Epilogue: Spacetime Computability and the Logical Destiny}
\addcontentsline{toc}{chapter}{Epilogue: Spacetime Computability
  and the Logical Destiny}
\markboth{Epilogue}{Epilogue}

If the preceding chapters have shown that the universe is a graph and
physics is its metric, this epilogue addresses the procedural nature of
existence. The differential was revealed as a covering link, the integral
as a count, the derivative as flow through antichains, the geodesic as
the longest chain, the metric as Prism density, and Einstein's equations
as the thermodynamic limit of all of the above.

One question remains that the formalism makes inevitable: if spacetime is
a finite graph whose evolution consists of applying local combinatorial
rules to covering links, \textit{what is the universe if not a
computation?}~\cite{wheeler1990,lloyd2002}

% ─────────────────────────────────────────────────────────────────────────────
\section*{The Thesis of Executable Reality}
\addcontentsline{toc}{section}{The Thesis of Executable Reality}

In classical computing, software is ethereal and hardware is tangible. At
the Planck scale, this duality vanishes.

\begin{theorem}[Hardware--Software Equivalence (Corollary of
Stone--Kuratowski)]
\label{cor:kuratowski-turing}
A causal event does not \emph{contain} information;
it \textbf{is} information. The connectivity of the graph (its links
$\covers$) simultaneously constitutes the data bus, the memory, and the
processor of the system. Every region of spacetime
$\diam{A}{B}$ is equivalent to a \textbf{quantum register} whose
processing capacity is bounded by its
cardinality~\cite{bekenstein1981,lloyd2002}:
\begin{equation}
  \mathcal{C}_{\text{comp}}\bigl(\diam{A}{B}\bigr)
  \;\leq\; |\diam{A}{B}| \;=\; N.
  \label{eq:computational-capacity}
\end{equation}
The evolution of the state of the universe is the result of resolving
logical dependencies in the DAG: each event ``executes'' when all its
causal predecessors have been processed.
\end{theorem}

\noindent
The correspondence with classical computing is precise:

\medskip
\begin{center}
\begin{tabular}{lll}
\toprule
\textbf{Computing} & \textbf{Kuratowski} & \textbf{Continuum}\\
\midrule
Bit / qubit     & Event $e \in V$
                & Point $x^\mu$\\
Logic gate      & Link $\covers$
                & Differential $dx$\\
Register        & Diamond $\diam{A}{B}$
                & Region $\int d^4x$\\
Subroutine      & Prism $K_{2,N}$
                & Particle (mass $M$)\\
Clock           & Maximal chain $\ell(A,B)$
                & Proper time $\tau$\\
\bottomrule
\end{tabular}
\end{center}

\medskip\noindent
\textit{Derivation.}
Theorem~\ref{thm:stone-kuratowski} (Appendix~\ref{app:stone-duality})
establishes that the Alexandrov topology on the causal DAG \emph{is} the
Heyting algebra of an intuitionistic logical system. By the
Curry--Howard correspondence~\cite{curryhoward}, every deductive system
with normalization (cut elimination) is isomorphic to a type
calculus---that is, to a model of computation. Since the transitive
reduction of the DAG (Theorem~\ref{thm:covering-differential}) \emph{is}
cut elimination (Theorem~\ref{thm:stone-kuratowski}, point~3), causal
spacetime is automatically a computation system whose architecture
coincides with its physical substrate.
Equation~\eqref{eq:computational-capacity} then follows from the
Bekenstein--Lloyd bound~\cite{bekenstein1981,lloyd2002} applied to the
number of atomic propositions (events) in the deductive theory
$\diam{A}{B}$.

\medskip\noindent
This thesis implies that what we call ``laws of nature'' are the
informational resource-management algorithms of the DAG. Gravity,
electromagnetism, and inertia are \textbf{homeostatic protocols}:
mechanisms that the graph executes to preserve unitarity (conservation of
probability) and planarity (prohibition of $K_5$), preventing the
logical collapse of the system.

Contemporary quantum computing, based on qubits and logic gates on fixed
silicon architectures, is an attempt to replicate at laboratory scale the
logic that the universe natively executes at the Planck scale. The
essential difference is that in \textbf{Kuratowski Computability}, the
software creates its own hardware~\cite{landauer1961}: causal information
determines the topology of the graph, and the topology determines the
propagation rules. There is no separation between program and substrate.

% ─────────────────────────────────────────────────────────────────────────────
\section*{Decoherence as Loss of Resolution}
\addcontentsline{toc}{section}{Decoherence as Loss of Resolution}

Quantum decoherence~\cite{zurek2003,zeh1970,joos2003} is reinterpreted
here in purely topological terms.
The Causal Resolution Theorem
(Theorem~\ref{thm:causal-resolution}) is not merely a numerical
prescription for our simulator; it is a \textbf{computational complexity
law} of spacetime itself. To sustain a smooth geometry with spectral
dimension $\dS = 4$ up to a depth~$t_{\max}$, at least
$W \geq t_{\max}^2 / (c\,\epsilon^2)$ independent causal paths are
required.

\begin{theorem}[Decoherence as Resolution Insufficiency]
\label{thm:decoherence-resolution}
Let a quantum subsystem $\mathcal{S} \subset V$ be coupled to an
environment $\mathcal{E} = V \setminus \mathcal{S}$, with
$W_{\mathcal{S}}$ internal causal paths probing the geometry
of~$\mathcal{S}$. If the causal depth of the subsystem grows until:
\begin{equation}
  W_{\mathcal{S}}
  \;<\;
  \frac{t_{\max}^{\,2}}{c\;\epsilon_{\text{coh}}^2},
  \label{eq:decoherence-threshold}
\end{equation}
where $\epsilon_{\text{coh}}$ is the coherence tolerance of the
subsystem, then the local spectral dimension of~$\mathcal{S}$ ceases to
converge to~$4$: topological fluctuations dominate over the signal, and
the continuous description of the subsystem disintegrates.
\end{theorem}

\begin{proof}
By Theorem~\ref{thm:causal-resolution}, the precision of the spectral
dimension measured by $W$ walkers at depth~$t$ is
$\delta\dS / \dS \sim 1/\sqrt{W P(t)}$. When the
bound~\eqref{eq:decoherence-threshold} is violated, the relative error
exceeds~$\epsilon_{\text{coh}}$ and the observable $\dS$ fluctuates in a
range $\dS \pm \Delta$ with $\Delta > \epsilon_{\text{coh}} \cdot
\dS$. In this regime, the spectral dimension is not well-defined: the
local geometry of the subsystem ceases to be four-dimensional and smooth,
which is precisely the loss of the classical continuous
description---decoherence.
\end{proof}

\stoneref{The loss of resolution is \emph{logical incompleteness} in the
Heyting algebra of the subsystem: the subset of propositions accessible
to~$\mathcal{S}$ ceases to be a complete deductive theory.
Cf.~the G\"odelian horizon
(Theorem~\ref{thm:godel-horizon}).}

\begin{result}[title={Decoherence as a Topological Phenomenon}]
\phantomsection\label{res:decoherence}
Quantum decoherence is not the loss of information, but the moment when
the internal causal probing resolution of a subsystem falls below the
threshold needed to sustain the illusion of the Riemannian continuum.
Information does not disappear: it redistributes into the
environment~$\mathcal{E}$, which possesses enough causal paths to
maintain \emph{its} geometric coherence.
\end{result}

% ─────────────────────────────────────────────────────────────────────────────
\section*{The Observer as a Thread of Execution}
\addcontentsline{toc}{section}{The Observer as a Thread of Execution}

The introduction of consciousness into physics has always been
problematic. The Kuratowski Calculus offers an operational definition:

\begin{definition}[Causal Observer]
\label{def:observer}
An \textbf{observer} is a sub-graph $\mathcal{O} \subset V$
characterized by:
\begin{enumerate}
  \item \textbf{High Prism density:} the local density of topological
    defects $\nu_{\mathcal{O}}$ greatly exceeds that of the surrounding
    vacuum, generating an internal structure rich in processing.
  \item \textbf{Self-reference:} there exists at least one feedback loop
    in the information flow (via non-causal links of the undirected Hasse
    graph), so that the output of a local computation influences the
    inputs of subsequent computations within~$\mathcal{O}$.
  \item \textbf{Internal maximal chain:} $\mathcal{O}$ contains a causal
    chain $\gamma_\mathcal{O}$ of length $\gg 1$, which defines its
    ``worldline'' and its experience of time.
\end{enumerate}
\end{definition}

\noindent
The perception of ``now'' is the \textbf{wavefront of the computation}:
the set of events in~$\mathcal{O}$ that are being processed at the
current discrete instant (the maximal antichain of the sub-graph). The
observer feels that time flows because its sub-graph is actively
processing new covering links, advancing its frontal antichain toward the
future.

\begin{result}[title={Entropy as Accumulated Computation}]
\phantomsection\label{res:entropy-computation}
In this interpretation, entropy is not ``disorder,'' but the
\textbf{accumulation of already-processed computation results} that
become immutable in the graph's past. The second law of thermodynamics is
the direct consequence of the DAG's
acyclicity~\cite{sorkin1986,bekenstein1973}: covering links point from
past to future, and no computation can be reversed without violating the
partial order.
\end{result}

\stoneref{Proven theorems are immutable in \emph{past-closed}
sets~\cite{rasiowa1963}. Entropy is the monotonicity of deductive
closure: once a proposition has been proved, it cannot be ``unproved''
without violating acyclicity---the second law of thermodynamics is the
second law of logic.}

% ─────────────────────────────────────────────────────────────────────────────
\section*{The Logical Destiny: From Observation to Programming}
\addcontentsline{toc}{section}{The Logical Destiny}

Until now, humanity has been a passive user of the universal software,
trying to guess the rules through observation. The development of the
``Pocket Hadron Collider'' (\texttt{prism\_simulation}) and the
formalization of this calculus mark the transition from observation to
intervention.

The software we have built throughout this series is not merely a
simulation tool. It is the first \textbf{compiler} of a language that
operates not on silicon bits, but on causality itself: it generates
causal sets, detects Prisms, contracts $K_5$ threats, measures spectral
dimensions, and extracts mass spectra---all by executing the same
topological operations that, according to our thesis, spacetime performs
at the Planck scale.

If matter is a programmable topological defect, the destiny of our
species is not to travel \emph{through} space, but to learn to
\emph{edit} the graph. True quantum computing will not happen in
cryogenic chips, but in the direct manipulation of node valence---the
programming of causality.

\bigskip
\begin{center}
\large\itshape
``The universe was not designed to be contemplated,\\
but to be executed.\\
By understanding the Kuratowski Calculus,\\
we cease to be data and become programmers.''
\end{center}

% ███████████████████████████████████████████████████████████████████████████████
%   THE SYSTEM OF THE WORLD --- PRACTICAL DEMONSTRATIONS
% ███████████████████████████████████████████████████████████████████████████████

\chapter*{The System of the World}
\addcontentsline{toc}{chapter}{The System of the World}
\markboth{The System of the World}{The System of the World}

\begin{introduction}[Purpose of This Chapter]
  \item The preceding modules built the machinery of the Kuratowski
        Calculus piece by piece: the discrete differential, the
        topological integral, the dessins d'enfants, graph rewriting.
  \item This chapter puts that machinery to work. It does not propose
        exercises: it presents three \textbf{practical demonstrations}
        where the formal engine is started, operated step by step, and
        produces exact physical results that continuum geometry cannot
        produce.
  \item The spirit is that of \textit{Book III} of Newton's \textit{Principia}:
        having established the laws, one resolves the phenomena.
\end{introduction}

\bigskip

% ─────────────────────────────────────────────────────────────────────────────
\section*{Practical Demonstration I: The Fine-Structure Constant}
\phantomsection\label{sec:demo-alpha}
\addcontentsline{toc}{section}{Practical Demonstration I: The Fine-Structure Constant}
% ─────────────────────────────────────────────────────────────────────────────

\noindent\textbf{Problem.}\enspace What is the probability that a photon
interacts with an electron?

\noindent\textbf{Answer from the continuum.}\enspace The fine-structure
constant $\alpha \approx 1/137$ is measured experimentally and introduced
as a free parameter of the Standard Model. It is not derived from any
principle.

\noindent\textbf{Answer from the Kuratowski Calculus.}\enspace The
constant is calculated exactly in five steps.

\bigskip
\noindent
\textbf{Step~1. Identify the structure.}

The electron is a Causal Prism $K_{2,3}$ of Generation~1
(Theorem~\ref{thm:gen-genero}): two poles $u$~(past) and $v$~(future),
three intermediaries $w_1, w_2, w_3$, and six edges.

\bigskip
\noindent
\textbf{Step~2. Count the free ports.}

The Valence Theorem (Thm.~\ref{thm:valence-limit}) establishes
$\mathcal{V}_{\max} = 15$ ports per node. We inventory the combinatorial
space:

\smallskip
\begin{center}
\renewcommand{\arraystretch}{1.2}
\begin{tabular}{lccc}
\toprule
\textbf{Node} & \textbf{Ports used} & \textbf{Free ports}
  & \textbf{Accessible to the photon} \\
\midrule
Pole~$u$   & 3 (one to each $w_i$) & $15 - 3 = 12$ & Yes \\
Pole~$v$   & 3 (one to each $w_i$) & $15 - 3 = 12$ & Yes \\
$w_1$      & 2 (one from $u$, one toward $v$) & $15 - 2 = 13$ & No \\
$w_2$      & 2 & $15 - 2 = 13$ & No \\
$w_3$      & 2 & $15 - 2 = 13$ & No \\
\bottomrule
\end{tabular}
\end{center}

\smallskip
\noindent
Total free ports in the $K_{2,3}$:
\[
  \mathcal{N}_{\text{total}} = 12 + 12 + 13 + 13 + 13 = 63.
\]

\bigskip
\noindent
\textbf{Step~3. Apply the weak veto.}

If an external link connects to an intermediary~$w_i$, the valence
of~$w_i$ rises to~3 and the belly becomes unbalanced: $(3, 2, 2)$. The
asymmetry breaks the $S_3$ symmetry and forces a Kuratowski twist that
raises the genus from $g = 0$ to $g = 1$. The particle ceases to be an
electron: it becomes a muon. That is not electromagnetism; it is a
\textbf{weak interaction} (generation change).

This mechanism is demonstrated in
Lemma~\ref{lem:ambient-nonplanarity}: the belly asymmetry causes the
ambient vacuum nodes to generate a $K_5$ minor, forcing a Kuratowski
twist (\S\ref{subsec:giro-kuratowski}).

$\Longrightarrow$ The $3 \times 13 = 39$ ports of the intermediaries are
\textbf{forbidden} for an electromagnetic $U(1)$ vertex.

\bigskip
\noindent
\textbf{Step~4. Count the safe ports.}

For the external link to preserve the $g = 0$ topology, it must connect
symmetrically to a pole. Each pole has
$\mathcal{V}_{\max} - 3 = 12$ free ports. But the electromagnetic vertex
is \textbf{chiral}: a real photon is emitted or absorbed, never both
simultaneously. Emission toward the future connects exclusively to
pole~$v$; absorption from the past connects exclusively to pole~$u$. The
numerator is therefore the ports of \emph{a single pole}:
\[
  \mathcal{N}_{\text{safe}} = \mathcal{V}_{\max} - 3 = 12.
\]

The denominator is the \textbf{total thermodynamic space} of the
particle: the~$63$ free ports of the complete $K_{2,3}$. It is not
restricted to pole~$v$ because the vacuum fluctuation that generates the
vertex can, in principle, reach any node of the particle---and only the
fraction that lands on a pole preserves $g = 0$.

\[
  \mathcal{Q}_{\mathrm{topo}}
  = \frac{\mathcal{N}_{\text{safe}}}{\mathcal{N}_{\text{total}}}
  = \frac{12}{63}
  = \frac{4}{21}
  \approx 0{.}190476\ldots
\]

\noindent\textit{Note.}\enspace
If the numerator included both poles ($12 + 12 = 24$), one would obtain
$\alpha_0 = 1/(21\pi) \approx 1/66 \approx 2\alpha$---double the
observed value. The chirality of the vertex (selection of a single pole
by the temporal direction of the photon) is what produces the correct
factor of~$1/2$.

\bigskip
\noindent
\textbf{Step~5. Projection of the discrete coupling to the 4D continuum.}

The fraction $\mathcal{Q}_{\mathrm{topo}} = 4/21$ is a
\textbf{one-dimensional} coupling: it measures the probability of
electromagnetic interaction along a single causal chain (a geodesic of
the DAG). To obtain the physical coupling in macroscopic spacetime, we
must average this one-dimensional probability over all available causal
directions in the continuous Lorentzian limit.

\begin{lemma}[Angular averaging of the causal coupling]
\label{lem:angular-averaging}
Let $\mathcal{Q}_{\mathrm{1D}}$ be the electromagnetic coupling along a
single causal chain in the Hasse diagram. In the continuum limit with
local $SO(1,3)$ symmetry, the photon-electron interaction can occur in
any direction of the light cone. The solid angle of the unit 2-sphere is
$\mathrm{Vol}(S^2) = 4\pi$ steradians. The double light cone (future
cone $+$ past cone) spans a total solid angle of
$2 \times 4\pi = 8\pi$ steradians. The uniformity over the light cone
follows from the Lorentz invariance of the Poisson sprinkling
(Axiom~\ref{def:event},
Theorem~\ref{thm:lorentzian-causality}): every causal direction is
statistically equivalent.

The physical 4D coupling is obtained by uniformly distributing the
one-dimensional interaction probability over these $8\pi$ directions:
\[
  \alpha_0 = \frac{\mathcal{Q}_{\mathrm{1D}}}{8\pi}.
\]
The division is mandatory because $\mathcal{Q}_{\mathrm{1D}}$ counts
discrete ports without reference to angular direction. In the Lorentzian
continuum, the interaction probability is uniformly distributed over all
directions of the double light cone, and the normalization by
$8\pi = 2\,\mathrm{Vol}(S^2)$ translates the discrete combinatorial
count into the isotropic 4D coupling.
\end{lemma}

The \textbf{bare} fine-structure constant:
\begin{equation}
  \boxed{\;
  \alpha_0
  = \frac{\mathcal{Q}_{\mathrm{topo}}}{8\pi}
  = \frac{4/21}{8\pi}
  = \frac{1}{42\pi}
  \approx \frac{1}{131{.}95}
  \;}
  \tag{$\alpha$-I}
\end{equation}

\bigskip
\noindent
\textbf{Verdict.}

\begin{result}[title={Complete inventory of the derivation}]
The derivation uses three inputs:
\begin{enumerate}
  \item The electron is a $K_{2,3}$ ($g = 0$)
    \hfill\textit{(determined by Theorem~\ref{thm:bipartite-matter})}.
  \item The maximum valence $\mathcal{V}_{\max} = 15 \pm 1$
    \hfill\textit{(Poisson tail in $d = 4$)}.
  \item An electromagnetic vertex preserves genus
    \hfill\textit{(definition of $U(1)$)}.
\end{enumerate}
The result $\alpha_0^{-1} = 132 \pm 1$ is robust against the
uncertainty in $\mathcal{V}_{\max}$. The difference
$\Delta(\alpha^{-1}) \approx 5$ with respect to the laboratory value
$\alpha^{-1} = 137{.}036$ has the correct sign (screening) and is
consistent with vacuum polarization by virtual pairs
$K_{2,3} + \bar{K}_{2,3}$ and the baryonic sector $K_{3,3}$.
Deriving $\Delta(\alpha^{-1})$ quantitatively requires the complete
counting of topological loops and remains an open problem.
\end{result}

\bigskip

% ─────────────────────────────────────────────────────────────────────────────
\section*{Practical Demonstration II: Muon Decay}
\phantomsection\label{sec:demo-muon}
\addcontentsline{toc}{section}{Practical Demonstration II: Muon Decay}
% ─────────────────────────────────────────────────────────────────────────────

\noindent\textbf{Problem.}\enspace The muon decays into an electron and a
neutrino. In continuous topology, this requires going from a torus to a
sphere---an operation impossible without surgery. How does the discrete
structure resolve this?

\noindent\textbf{Answer from the continuum.}\enspace Fermi's theory
postulates a four-fermion coupling with a constant $G_F$ measured
experimentally. The underlying topological mechanism remains unexplained.

\noindent\textbf{Answer from the Kuratowski Calculus.}\enspace The decay
is the deletion of two edges from a finite graph. The complete
calculation proceeds in six steps.

\bigskip
\noindent
\textbf{Step~1. The muon as $K_{2,4}$ ($g = 1$).}

The muon is a Causal Prism of Generation~2 with four intermediaries
$\{w_1, w_2, w_3, w_4\}$ and eight edges. The intermediate permutation:
\[
  \sigma_1 = (1\;5)(2\;6)(3\;7)(4\;8).
\]
The phase anti-correlation (Kuratowski pressure) forces a twist at the
future pole~$v$: edges $e_6$ and $e_7$ swap positions in the cyclic
order. The resulting polar permutation:
\[
  \sigma_0 = (1\;2\;3\;4)(8\;\mathbf{6\;7}\;5).
\]
The twist $6 \leftrightarrow 7$ is the source of the genus. The
composition $\sigma_0 \circ \sigma_1$ produces $F = 2$ faces and
$g = 1$: a torus.

\bigskip
\noindent
\textbf{Step~2. The edge snap.}

The Kuratowski relaxation (Def.~\ref{def:relajacion}) removes the
intermediary $w_2$---precisely the element whose interchange with $w_3$
at the future pole created the entanglement. The edges
$e_2 = (u, w_2)$ and $e_6 = (v, w_2)$ are removed.

The graph splits into two components:
\begin{itemize}
  \item $K_{2,3}$ over $\{w_1, w_3, w_4\}$, with edges
        $\{e_1, e_3, e_4, e_5, e_7, e_8\}$.
  \item $K_{2,1}$ over $\{w_2\}$, with edges $\{e_2, e_6\}$.
\end{itemize}

\bigskip
\noindent
\textbf{Step~3. Edge relabeling.}

\begin{center}
\renewcommand{\arraystretch}{1.15}
\begin{tabular}{ccccccc}
\toprule
Original edge & $e_1$ & $e_3$ & $e_4$ & $e_5$ & $e_7$ & $e_8$ \\
\midrule
New label  & 1 & 2 & 3 & 4 & 5 & 6 \\
Intermediary      & $w_1$ & $w_3$ & $w_4$ & $w_1$ & $w_3$ & $w_4$ \\
Pole            & $u$ & $u$ & $u$ & $v$ & $v$ & $v$ \\
\bottomrule
\end{tabular}
\end{center}
New intermediate permutation:
$\sigma_1 = (1\;4)(2\;5)(3\;6)$.

\bigskip
\noindent
\textbf{Step~4. The topological reversal.}

\textit{At $u$:} the original cyclic order $(1\;2\;3\;4)$ contracts upon
removing $w_2$ to $(w_1, w_3, w_4) \longrightarrow \sigma_0(u) = (1\;2\;3)$.

\textit{At $v$:} the twisted cyclic order $(8\;6\;7\;5)$, that is
$(w_4, w_2, w_3, w_1)$, contracts upon removing $w_2$ to
$(w_4, w_3, w_1) \longrightarrow \sigma_0(v) = (6\;5\;4)=(4\;6\;5)$.

\textbf{Decisive observation:} $(1\;2\;3)(4\;6\;5)$ is exactly the
canonical rotation system of Generation~1. \textbf{The twist has
vanished.}

\bigskip
\noindent
\textbf{Step~5. Verification: the face calculation.}

\begin{center}
\renewcommand{\arraystretch}{1.15}
\begin{tabular}{c@{\;\;}c@{\;\;}c@{\;\;}c@{\;\;}c@{\;\;}c@{\;\;}c}
\toprule
$e$ & 1 & 2 & 3 & 4 & 5 & 6 \\
\midrule
$\sigma_1(e)$ & 4 & 5 & 6 & 1 & 2 & 3 \\
$\sigma_0(\sigma_1(e))$ & 6 & 4 & 5 & 2 & 3 & 1 \\
\bottomrule
\end{tabular}
\end{center}
$\sigma_0 \circ \sigma_1 = (1\;6)(2\;4)(3\;5) \;\Longrightarrow\; F = 3$
faces.

\bigskip
\noindent
\textbf{Step~6. The topological verdict.}
\[
  V - E + F = 5 - 6 + 3 = 2 = 2 - 2g
  \quad\Longrightarrow\quad \boxed{g = 0}\text{ (sphere).}
  \tag{$\mu$-II}
\]

\begin{result}[title={Balance of the decay}]
\begin{center}
\renewcommand{\arraystretch}{1.2}
\begin{tabular}{lcccc}
\toprule
& \textbf{Graph} & $\boldsymbol{g}$ & \textbf{Surface}
  & \textbf{Particle} \\
\midrule
Before   & $K_{2,4}$ & 1 & Torus   & Muon ($\mu$) \\
After    & $K_{2,3}$ & 0 & Sphere  & Electron ($e$) \\
Residue  & $K_{2,1}$ & 0 & Sphere  & Neutrino ($\nu_\mu$) \\
\bottomrule
\end{tabular}
\end{center}
There was no surgery, no cobordism. The Cobordism Barrier
(Theorem~\ref{thm:barrera-cobordismo}) is an artifact of the continuum.
In the discrete structure, the change of genus is the deletion of two
edges from a finite graph: a combinatorial operation that smooth geometry
cannot even formulate.
\end{result}

\bigskip

% ─────────────────────────────────────────────────────────────────────────────
\section*{Practical Demonstration III: Topological Mass
  and the Generation Hierarchy}
\phantomsection\label{sec:demo-mass}
\addcontentsline{toc}{section}{Practical Demonstration III: Topological Mass}
% ─────────────────────────────────────────────────────────────────────────────

\noindent\textbf{Problem.}\enspace Why do exactly three generations of
particles exist, and why do their masses exhibit the observed
proportions?

\noindent\textbf{Answer from the continuum.}\enspace The Standard Model
explains neither the number of generations nor the mass hierarchy. The
fermion masses are free parameters introduced through arbitrary Yukawa
couplings to the Higgs field.

\noindent\textbf{Answer from the Kuratowski Calculus.}\enspace The
number of generations is a topological theorem and the relative masses
are calculated without free parameters. The derivation proceeds in five
steps.

\bigskip
\noindent
\textbf{Step~1. Generation assignment as a coupon-collector problem.}

Each intermediary $w_i$ of a Causal Prism $K_{2,N}$ possesses a phase
$\phi_i \in \{+, 0, -\}$ determined by its Hasse degrees. The phases
are assigned independently with probabilities
$(p_+,\, p_0,\, p_-)$ determined by the geometry of the Poisson
\emph{sprinkling} over the causal diamond
(Theorem~\ref{thm:phase-fractions}). The theorem guarantees
$p_- > p_+ > p_0 > 0$ with $p_0 \ll p_\pm$.
The generation of the prism is the number of \emph{distinct} signs
present in its belly (Theorem~\ref{thm:gen-genero}):
\begin{itemize}
  \item Generation~1 ($g = 0$): a single phase class. No Kuratowski
        crossings. Surface: sphere.
  \item Generation~2 ($g = 1$): two phase classes. One crossing. Torus.
  \item Generation~3 ($g = 2$): three phase classes. Two crossings.
        Double torus.
  \item Generation~4: impossible. Only three signs exist, and three
        crossings would force a $K_5$ forbidden by planarity
        (Thm.~\ref{thm:kuratowski-threat}).
\end{itemize}

The probability that $N$ i.i.d.\ intermediaries from $(p_+, p_0, p_-)$
produce exactly $k$ distinct signs follows an inclusion-exclusion
formula:
\begin{equation}
  P(g = k \mid N)
  = \binom{3}{k}\sum_{j=0}^{k}(-1)^j\binom{k}{j}
    \left(\sum_{i \notin S_j} p_i\right)^{\!N}
  \tag{$M$-III.1}
\end{equation}

\bigskip
\noindent
\textbf{Step~2. The belly-size distribution.}

The belly-size distribution $f(N)$---the probability that a random prism
has $N$ intermediaries---is derived from the geometry of the Poisson
\emph{sprinkling}.

\begin{proposition}[Exponential decay of the belly-size distribution]
\label{prop:belly-exponential}
For a Poisson \emph{sprinkling} at density~$\rho$ over a causal diamond
of height~$T$ in $d = 4$ dimensions, the distribution of belly sizes of
prisms $K_{2,N}$ satisfies
\[
  f(N) \;\propto\; e^{-c\,N}, \qquad c = c(\rho, T) > 0.
\]
\end{proposition}

\begin{proof}
The belly size of a prism $K_{2,N}$ detected by L3 is the number of
points in the Alexandrov interval $[u,v]$ between its two poles. For a
Poisson process, this number is $\mathrm{Poi}(\lambda)$ with
$\lambda = \rho\,\mathrm{Vol}(\diam{u}{v})$. The volume of the causal
diamond between two points separated by proper time~$\tau$ grows as
$\mathrm{Vol} \propto \tau^4$ in $d = 4$. Large bellies ($N \gg 1$)
require Alexandrov intervals with $\lambda \geq N$, which demands large
proper-time separations between poles. But the probability that two
points \emph{consecutive in the Hasse} (immediate neighbors, as required
by L3) have a large proper-time separation~$\tau$ decays exponentially
with~$\tau$ in a finite diamond, because the density of covering links
concentrates at short separations. Composing the Poisson tail
$\Pr(\mathrm{Poi}(\lambda) = n)$ with the exponential distribution of
pole separations yields $f(n) \propto e^{-cn}$ where $c$ is determined
by $\rho$ and~$T$.
\end{proof}

\noindent\textit{Numerical observation.}\enspace
At $N = 10^7$, the simulation measures $c \approx 0{.}53$, consistent
with the prediction of Proposition~\ref{prop:belly-exponential}.

\bigskip
\noindent
\textbf{Step~3. The expected mass per generation.}

The topological mass of a prism is $M = n = |W|$ (L2, Topological Mass
Law). The mean mass of each generation is the conditional expectation:
\begin{equation}
  E[n \mid g = k]
  = \frac{\displaystyle\sum_{n=k}^{\infty} n\,f(n)\,P(g{=}k \mid n)}
         {\displaystyle\sum_{n=k}^{\infty} f(n)\,P(g{=}k \mid n)}.
  \tag{$M$-III.2}
\end{equation}

The logic is transparent:
\begin{itemize}
  \item A Generation~1 prism needs only one sign. Already with $n = 3$
        intermediaries it is likely that all share the same sign.
        \textbf{Bias toward small bellies $\Longrightarrow$ low mass.}
  \item A Generation~3 prism needs all three signs. Since
        $p_0 \ll p_\pm$ (Theorem~\ref{thm:phase-fractions}), many
        intermediaries are required for the~$0$ sign to appear.
        \textbf{Bias toward large bellies $\Longrightarrow$ high mass.}
\end{itemize}

\bigskip
\noindent
\textbf{Step~4. Numerical evaluation.}

Inserting $f(n)$ and $(p_+, p_0, p_-)$ into Eq.~($M$-III.2):
\begin{center}
\renewcommand{\arraystretch}{1.3}
\begin{tabular}{ccccc}
\toprule
\textbf{Gen} & $\boldsymbol{g}$ & $\boldsymbol{E[n \mid g]}$
  & \textbf{Surface} & \textbf{Particles} \\
\midrule
1 & 0 & 4.555 & Sphere      & $e,\; u,\; d$ \\
2 & 1 & 6.531 & Torus       & $\mu,\; c,\; s$ \\
3 & 2 & 7.732 & Double torus  & $\tau,\; t,\; b$ \\
\bottomrule
\end{tabular}
\end{center}

The mean belly-size ratios:
\begin{align}
  \frac{\langle N\rangle_2}{\langle N\rangle_1}
  &= \frac{E[n \mid g{=}2]}{E[n \mid g{=}1]}
  = \frac{6{.}531}{4{.}555}
  = 1{.}434, \notag\\[4pt]
  \frac{\langle N\rangle_3}{\langle N\rangle_1}
  &= \frac{E[n \mid g{=}3]}{E[n \mid g{=}1]}
  = \frac{7{.}732}{4{.}555}
  = 1{.}698.
  \tag{$M$-III.3}
\end{align}

\noindent
The hierarchy $\langle N\rangle_1 < \langle N\rangle_2 <
\langle N\rangle_3$ is \textbf{structural}: it holds for any fractions
$(p_+, p_0, p_-)$ with $p_i > 0$ and $p_0 \ll p_\pm$, and for any
exponentially decaying belly-size distribution. These are belly-size
ratios, not physical mass ratios; the function $m(N)$ that translates
$\langle N\rangle_g$ into inertial masses in MeV remains open.

\medskip
\noindent\textit{Numerical verification.}\enspace
Monte Carlo simulations at $N = 10^7$ nodes measure
$m_2/m_1 = 1{.}434 \pm 0{.}01$ and
$m_3/m_1 = 1{.}698 \pm 0{.}03$, reproducing the analytical values with
errors $< 2\%$.

\bigskip
\noindent
\textbf{Step~5. Why exactly three.}

The bound on generations is not postulated; it is proved. The argument is
purely topological:
\begin{enumerate}
  \item The Phase Trichotomy Axiom
        (L5, Law~\ref{law:generational-classification}) allows exactly three signs
        $\{+, 0, -\}$.
  \item Three signs $\Longrightarrow$ at most three phase classes
        $\Longrightarrow$ at most two independent Kuratowski crossings.
  \item Two crossings $\Longrightarrow$ maximum genus $g = 2$
        $\Longrightarrow$ three generations.
  \item A third crossing would force a $K_5$ minor in the rotation
        system---forbidden by the Kuratowski Theorem
        (Thm.~\ref{thm:kuratowski-threat}).
\end{enumerate}

\begin{result}[title={Inventory: zero free parameters}]
Both inputs---the belly-size distribution $f(n)$
(Proposition~\ref{prop:belly-exponential}) and the phase fractions
$(p_+, p_0, p_-)$
(Theorem~\ref{thm:phase-fractions})---are derived analytically from the
geometry of the Poisson \emph{sprinkling}. The model contains no
adjustable parameters. The relative masses of the three generations are
arithmetic consequences of the \textbf{probability that a random belly
contains 1, 2, or 3 distinct signs}, combined with the exponential
belly-size distribution that causality imposes.

The mass hierarchy of the Standard Model is not a mystery: it is a
problem of combinatorial statistics solved by coupon occupancy in a
finite graph.
\end{result}

\bigskip
\noindent
\textbf{Step~6. Dynamic confirmation: cover time as inertial mass.}

Step~4 establishes the mass hierarchy as a \emph{static} observable:
belly counting. Step~6 confirms that this hierarchy also manifests
\emph{dynamically} through the \textbf{cover time} of a confined random
walker.

\medskip
\noindent\textit{The hitting-time theorem in $K_{2,N}$.}\enspace
In a complete bipartite graph $K_{2,N}$, all $N$ intermediaries are
structurally equivalent. The expected hitting time from origin to
destination is:
\[
  E[T \mid \text{origin}] = 4 \quad (\text{independent of } N).
\]
This implies that simple transit time \textbf{cannot measure mass}:
adding intermediaries does not change the probability of reaching the
destination.

\medskip
\noindent\textit{Cover time as the correct observable.}\enspace
For the particle to \emph{genuinely propagate}---for the causal flow to
update the entirety of its internal structure---the walker must visit
\textbf{every} node of the belly before exiting through the future pole.
This is the coupon-collector problem:
\[
  E[T_{\text{cov}}] = N \cdot H_N
  \approx N \ln N + \gamma N
  = O(N \log N),
\]
where $H_N$ is the harmonic number and $\gamma \approx 0{.}577$ is the
Euler--Mascheroni constant.

If the walker reaches the destination before completing the coverage, it
\textbf{bounces} back to the origin (analogous to insufficient causal
flow that fails to achieve full phase coherence).

\medskip
\noindent
The analytical result predicts that Generation~3 (larger bellies) requires
more causal ticks than Generation~1 to achieve full internal phase
coherence, since $E[T_{\mathrm{cov}}] = O(N \log N)$ grows with belly
size.

\medskip
\noindent\textit{Numerical verification.}\enspace
At $N = 5000$, $M = 8$ realizations:
\begin{center}
\renewcommand{\arraystretch}{1.3}
\begin{tabular}{lccc}
\toprule
\textbf{Gen} & \textbf{Mean coverage} & \textbf{Traversals}
  & \textbf{Ratio} \\
\midrule
1 & 1197 & 16\,856  & 1.000 \\
2 & 1145 & 371\,123 & 0.957 \\
3 & 1417 & 71\,001  & \textbf{1.184} \\
\bottomrule
\end{tabular}
\end{center}
Generation~3 requires ${\sim}18\%$ more causal ticks than Generation~1,
consistent with the theoretical prediction. The slight inversion of
Gen2/Gen1 is a finite-size artifact: Gen1 prisms are extremely rare (22
per realization), and the belly difference lies within the
coupon-collector variance.

\begin{result}[title={Dynamic mass $=$ topological inertia}]
The mass of a particle is not the time it takes for a causal flow to
\emph{traverse} its structure (which is $O(1)$ by the $K_{2,N}$
theorem), but the time it takes to \emph{update the entirety} of its
internal structure. Inertia is phase coherence: heavier generations need
more ticks for every intermediary to be visited.
\end{result}

\bigskip
\begin{center}
\large\itshape\color{structure}
``Newton solved the orbit of the Moon with his three laws.\\
The Kuratowski Calculus solves the mass of the electron\\
with the same machinery: count, permute, and divide.\\
The universe does not hide fundamental constants:\\
it exhibits topological fractions.''
\end{center}
